\section{Final design results}
\label{sec:final-results}
Table~\ref{tab:final-result} is generated from the current final-design model. Source-backed process quantities are not mixed with historical Gate-5 or Review-5 deployment values.

\begin{table}[htbp]\centering\scriptsize
\caption{Final-design CN4252 results and required economic sensitivities.}
\label{tab:final-result}
\resizebox{\linewidth}{!}{\pgfplotstabletypeset[col sep=comma,string type,
columns={scenario,h2ty,heatmw,heliumkgs,capturedty,avoidedty,incrementalsgd,costsgdt,apass,cpass,joint},
columns/scenario/.style={column name={Scenario}},columns/h2ty/.style={column name={H$_2$ t/y}},
columns/heatmw/.style={column name={Heat MWth}},columns/heliumkgs/.style={column name={He kg/s}},
columns/capturedty/.style={column name={CO$_2$ captured t/y}},columns/avoidedty/.style={column name={Avoided tCO$_2$e/y}},
columns/incrementalsgd/.style={column name={Incremental S\$/y}},columns/costsgdt/.style={column name={S\$/t}},
columns/apass/.style={column name={A}},columns/cpass/.style={column name={C}},columns/joint/.style={column name={Joint}},
every head row/.style={before row=\toprule,after row=\midrule},every last row/.style={after row=\bottomrule}]
{../results/final_design/generated/final_design_results.csv}}
\end{table}

The final source process produces approximately 98 ktH$_2$/y at 85\% availability. Direct source-balance abatement is already well above 0.25 Mt/y because the conventional no-capture baseline emits 3,205 short ton/day CO$_2$ while the nuclear-assisted captured case emits 142 short ton/day at the same H$_2$ output~\cite{inltev953,inltev961}. The lifecycle extension retains a substantial positive margin after upstream-gas, nuclear, auxiliary and transport/storage burdens.

\subsection{Temperature and architecture sensitivity}
The final reformer temperature remains 871$^\circ$C across the INL sensitivity because it is the selected process optimum. At 850$^\circ$C delivered heat (875$^\circ$C ROT), fired trim remains necessary; at 900$^\circ$C delivered heat (925$^\circ$C ROT), nuclear heat supplies the full reforming duty~\cite{inltev961}. This is the principal temperature sensitivity: lower reactor/process-heat temperatures do not change the target chemistry but change how much fossil trim is required.

\begin{table}[htbp]\centering\scriptsize
\caption{Final-design temperature architecture and relevant source sensitivity.}
\resizebox{.85\linewidth}{!}{\pgfplotstabletypeset[col sep=comma,string type,
columns={case,reactoroutc,heatdeliveryc,reformerc,architecture},
every head row/.style={before row=\toprule,after row=\midrule},every last row/.style={after row=\bottomrule}]
{../results/final_design/generated/temperature_sensitivity.csv}}
\end{table}

\subsection{CN4252 interpretation}
The final design is classified only after forward physical and economic calculation. A joint pass, if reproduced in Table~\ref{tab:final-result}, is a \textbf{CONDITIONAL MODEL RESULT}: it depends on the mature GTHTR300C design-study cost basis, an enlarged IHX/secondary-loop train, cross-border CCS, the stated lifecycle proxies and electricity-value assumptions. It is not evidence that a Singapore commercial project presently exists at those costs.
